\documentclass[11pt]{article}

\usepackage[margin=0.7in]{geometry}
\usepackage{amsmath}
\usepackage{amssymb}
\usepackage{booktabs}
\usepackage{array}
\usepackage{hyperref}
\usepackage{enumitem}
\usepackage{listings}
\usepackage{xcolor}
\usepackage{parskip}
\usepackage{titlesec}

\hypersetup{colorlinks=true, linkcolor=blue, urlcolor=blue}

\titlespacing*{\section}{0pt}{1.1em}{0.5em}
\titlespacing*{\subsection}{0pt}{0.9em}{0.4em}
\titlespacing*{\paragraph}{0pt}{0.7em}{0.5em}
\setlist{itemsep=2pt, topsep=2pt, parsep=0pt}

\lstset{
    basicstyle=\ttfamily\small,
    breaklines=true,
    frame=single,
    columns=fullflexible,
}

\setlength{\parskip}{0.5em}
\setlength{\parindent}{0pt}

\title{\vspace{-1.5em}Sequence Alignment: Basic DP vs.\ Hirschberg's Memory-Efficient Algorithm}
\author{}
\date{}

\begin{document}

\maketitle
\vspace{-3em}

This project implements and compares two algorithms for the global DNA sequence
alignment problem: the classical Needleman--Wunsch dynamic programming (DP)
algorithm and Hirschberg's linear-space divide-and-conquer algorithm. Both
algorithms always find an alignment of the same optimal cost under a fixed gap
penalty and mismatch cost matrix, but they differ sharply in memory use: the
basic DP algorithm needs $O(mn)$ space, while Hirschberg's algorithm reduces
this to $O(m+n)$ at the cost of a modest constant-factor slowdown in running
time.

\section{Problem Statement}

Given two DNA strings $X$ and $Y$ over the alphabet $\{A, C, G, T\}$, the
\emph{global sequence alignment} problem asks for the lowest-cost way to align
$X$ and $Y$ by inserting gap characters (`\_') into either string so that the
resulting strings have equal length. Each aligned column incurs a cost: a
\textbf{gap penalty} $\delta$ if one character in the column is a gap, or a
\textbf{mismatch/match cost} $\alpha(x,y)$ otherwise, taken from a fixed
substitution cost matrix. The goal is to minimize the total cost summed over
all columns.

\section{Scoring Scheme}

\textbf{Gap penalty:} $\delta = 30$

\textbf{Mismatch cost matrix} $\alpha(x,y)$:

\begin{center}
\begin{tabular}{c|cccc}
\toprule
 & A & C & G & T \\
\midrule
A & 0   & 110 & 48  & 94  \\
C & 110 & 0   & 118 & 48  \\
G & 48  & 118 & 0   & 110 \\
T & 94  & 48  & 110 & 0   \\
\bottomrule
\end{tabular}
\end{center}

\section{Input Format}

Rather than storing raw DNA strings, each input file stores a compressed
description of each sequence: an initial seed string over $\{A, C, G, T\}$,
plus a list of 0-indexed insertion positions. The full sequence is
reconstructed by repeatedly duplicating the current string and splicing the
copy in at the given position:
\[
\texttt{result} \leftarrow \texttt{result}[0:\texttt{pos}+1] \;+\;
\texttt{result} \;+\; \texttt{result}[\texttt{pos}+1:]
\]
This roughly doubles the string length at each step, so long test sequences
can be specified compactly. Both $X$ and $Y$ are encoded this way, one after
the other, in a single input file.

\section{Algorithms}

\subsection{Basic Dynamic Programming (\texttt{basic.py})}

The standard Needleman--Wunsch recurrence builds a full $(m{+}1) \times (n{+}1)$
cost table \texttt{OPT}, where $m = |X|$ and $n = |Y|$:
\[
OPT[i][j] =
\begin{cases}
  i \cdot \delta & j = 0 \\
  j \cdot \delta & i = 0 \\
  \min\big(OPT[i-1][j-1] + \alpha(X_i, Y_j),\; OPT[i-1][j] + \delta,\;
  OPT[i][j-1] + \delta\big) & \text{otherwise}
\end{cases}
\]
The optimal cost is $OPT[m][n]$; the alignment strings are recovered by
tracing back through the full table from $(m,n)$ to $(0,0)$.

\textbf{Complexity:} Time $O(mn)$, Space $O(mn)$ --- the whole table must stay
in memory to support traceback.

\subsection{Memory-Efficient Divide and Conquer (\texttt{efficient.py})}

This implementation uses \textbf{Hirschberg's algorithm}, which reaches the
same optimal cost in the same $O(mn)$ time but uses only $O(m+n)$ space, by
never materializing the full table.

\textbf{Linear-space cost rows.} \texttt{last\_row\_costs(X, Y)} computes only
the final row of the DP table, keeping just two rows (previous/current) in
memory at a time.

\textbf{Divide step.} Split $X$ at its midpoint into $X_{\text{left}}$ and
$X_{\text{right}}$; compute the forward cost row for $X_{\text{left}}$ vs.\
$Y$, and the backward cost row for reversed $X_{\text{right}}$ vs.\ reversed
$Y$. The column $j^*$ minimizing the sum of forward and backward cost marks
where the optimal path crosses the midpoint row.

\textbf{Conquer step.} Recurse on $(X_{\text{left}}, Y[0{:}j^*])$ and
$(X_{\text{right}}, Y[j^*{:}])$ and concatenate the results. The recursion
bottoms out at a small full DP + traceback (\texttt{small\_dp}) once either
sequence has length $\le 1$.

\textbf{Complexity:} Time $O(mn)$ (recomputed cost rows at each recursion
level still sum to $O(mn)$ overall), Space $O(m+n)$.

\section{Correctness Verification}

\texttt{compare.py} is a randomized differential test: it generates random
DNA strings (length 1--80), runs both algorithms, and checks they report the
\emph{same optimal cost}. It does not require byte-identical alignment
strings, since ties at the optimal cost can be broken differently by the two
algorithms. Over 200 trials: \textbf{identical optimal cost 200/200 (100\%)},
byte-identical alignment strings 48/200 (24\%) --- both expected outcomes,
since cost-optimality is guaranteed but the alignment itself need not be
unique.

\section{Benchmark Methodology}

\texttt{benchmark.py} sweeps problem size $m+n \in \{128, 256, 512, 1024,
2048, 3072, 4096\}$ and measures CPU time (\texttt{time.perf\_counter()}) and
peak memory via Python's \texttt{tracemalloc} (actual bytes allocated, a
cleaner signal than an RSS delta). Results are written to \texttt{results.csv}.

\section{Results}

\begin{center}
\begin{tabular}{r r r r r r r}
\toprule
$m+n$ & Basic (ms) & Efficient (ms) & Basic peak (KB) & Efficient peak (KB) & Mem.\ saved & Time cost \\
\midrule
128  & 4.0    & 6.6    & 147.4    & 14.1  & 90.4\% & 1.67$\times$ \\
256  & 13.9   & 23.5   & 588.7    & 19.0  & 96.8\% & 1.69$\times$ \\
512  & 47.7   & 89.1   & 2{,}337.4 & 37.9  & 98.4\% & 1.87$\times$ \\
1024 & 300.9  & 419.4  & 9{,}291.5 & 76.2  & 99.2\% & 1.39$\times$ \\
2048 & 1{,}455.9 & 1{,}914.2 & 37{,}025.7 & 150.8 & 99.6\% & 1.31$\times$ \\
3072 & 3{,}471.9 & 4{,}702.0 & 83{,}189.3 & 224.7 & 99.7\% & 1.35$\times$ \\
4096 & 6{,}422.1 & 8{,}697.2 & 147{,}785.5 & 298.5 & \textbf{99.8\%} & 1.35$\times$ \\
\bottomrule
\end{tabular}
\end{center}

\textbf{Time.} Both algorithms scale as $O(mn)$ -- time roughly quadruples
each time $m+n$ doubles. Hirschberg's is consistently \textbf{1.3--1.9$\times$
slower}, since the divide step recomputes forward/backward cost rows at every
recursion level; this is the expected cost of trading memory for arithmetic.

\textbf{Memory.} The basic DP table grows quadratically; Hirschberg's $O(m+n)$
rows grow linearly. The savings widen with input size, from 90.4\% at
$m+n=128$ to \textbf{99.8\%} at $m+n=4096$ (147.8 MB down to 0.3 MB).

\textbf{Takeaway.} Hirschberg's algorithm is the clear choice once sequence
length makes the quadratic-memory basic DP impractical, at the cost of a
modest, bounded runtime penalty.

\section{Project Structure}

\begin{lstlisting}
Sequence-Alignment/
|-- basic.py          Basic Dynamic Programming implementation
|-- efficient.py      Hirschberg's memory-efficient implementation
|-- basic.sh          Runs the basic implementation
|-- efficient.sh      Runs the efficient implementation
|-- compare.py        Differential test: both algorithms vs. the DP optimum
|-- benchmark.py      Time/memory sweep across input sizes -> results.csv
|-- results.csv       Measured output of benchmark.py
|-- sample_input.txt  Small example in the input format described above
|-- Summary.pages     Project report
\end{lstlisting}

\section{Running the Project}

Both commands below run on \texttt{sample\_input.txt} as-is and print an
optimal alignment cost of \textbf{660}.

\begin{lstlisting}
# Basic algorithm
./basic.sh sample_input.txt output.txt
python3 basic.py sample_input.txt output.txt

# Memory-efficient algorithm
./efficient.sh sample_input.txt output.txt
python3 efficient.py sample_input.txt output.txt

# Correctness check and benchmark
python3 compare.py
python3 benchmark.py
\end{lstlisting}

Each run writes an output file containing, in order: the minimum alignment
cost, the two aligned sequences, the execution time, and the peak memory
usage.

\section{Requirements}

Python 3.x and \texttt{psutil} (\texttt{pip install psutil}), used by
\texttt{basic.py} and \texttt{efficient.py} to report peak memory. The
alignment algorithms themselves use only the Python standard library;
\texttt{benchmark.py} and \texttt{compare.py} need nothing beyond that.

\section{Applications}

DNA sequence alignment, computational biology, bioinformatics, genome
comparison, evolutionary analysis.

\section{References}

\begin{itemize}
    \item Needleman, S.~B., \& Wunsch, C.~D. (1970). \textit{A General Method
    Applicable to the Search for Similarities in the Amino Acid Sequence of
    Two Proteins.} Journal of Molecular Biology, 48(3), 443--453.
    \item Hirschberg, D.~S. (1975). \textit{A Linear Space Algorithm for
    Computing Maximal Common Subsequences.} Communications of the ACM, 18(6),
    341--343.
\end{itemize}

\end{document}
